\scanpage{11}
\mspage{9}

\unsure{\deleted{L}C|\deleted{C}L}\unsure{\deleted{a}e|\deleted{e}a}s\deleted{\unsure{e|c}} bicomplexes s'identifient aux complexes de chhines de lo\unsure{\deleted{g}n|\deleted{n}g}gueur 2.
Donc la catégorie de (2-Cat)-groupes abéliens est équivelente à celle de
complexes de chaines de longueur 2. L'explicitation de la 2-catégorie
associée à un tel complexe de chaines est mon interprétation chérie. Il
faudrait encore expliciter les \deleted{morphismes} 2-Grab-catégorie des
homomorphismes (en un sens à préciser) entre deu\unsure{\deleted{t}x|\deleted{x}t}\deleted{\unsure{e}l} tel\unsure{\deleted{s}l|\deleted{l}s}es 2-Grab-catég
ries strictes, en termes d'un complexe RHom.

L'idée qui se dégage de cette approche, et que Quillen on espère v
expl\unsure{\deleted{o}i|\deleted{i}o}citer, c'est que la \deleted{notion de catégorie} de notion\deleted{s} de groupe \uline{abélie}
dans (n-Cat) correspond à celle de complexe de \inserted{\begin{tabular}[b]{@{}l@{}}chaines de\\[-.9ex]
\begin{tikzpicture}[x=1cm,y=1cm]\draw[line width=.45pt] (0,0)..controls(-.08,.12)and(-.02,.3)..(.20,.47)--(.24,.49);\draw[line width=.4pt] (.27,.50)..controls(.75,.61)and(1.19,.57)..(1.55,.55)--(1.70,.53)--(1.83,.54)..controls(2.05,.55)and(2.29,.57)..(2.38,.64)..controls(2.42,.68)and(2.43,.73)..(2.41,.77);\end{tikzpicture}
\end{tabular}} longueur n, en raisonnant
soit mod. isomorphisme de part et d'autre, soit mod. une notion naturell
de n-équivalence d'un côté, et dans la catégorie dérivée de l'autre. Cec
doit pouvoir se préciser en utilisant des Hom$^{\bullet}$ resp. des RHom$^{\bullet}$.

Si maintenant C est une catégorie avec un ``produit tensoriel'' assoc
tif (et éventuellement unitaire), i.e. un bifoncteur avec des isomorphis%
mes d'associativité (mais sans inverse ...), on peut lui associer \unsure{\deleted{(}s|\deleted{s}(}ans
doute une Gr-catégorie universelle (tout comme à un monoïde on associe
groupe symmétrisé), d'où des invariants, qu'on pourra désigner par K$_{o}$(C)
et K$_{1}$(C). On suppose qu'on pourra de même trouver, pour tout n, une
Gr-n-catégorie universelle où envoyer C de façon compatib\unsure{\deleted{o}l|\deleted{b}l|\deleted{l}o|\deleted{l}b}e avec $\otimes$ (\unsure{\deleted{d}t|\deleted{t}d}out
cela dans un sens à préciser), dont les invariants $\handwritten{\pi}_{n}$ seraient par déf
tion les K$_{n}$(C). \unsure{\deleted{E}U|\deleted{L}U|\deleted{t}U|\deleted{U}E|\deleted{U}L|\deleted{U}t}\unsure{\deleted{a}n|\deleted{n}a} cas intéressant est celui où C est une catégorie ave
produits finis (par exemple une catégorie additive) et $\otimes$ est le produit;
\deleted{mais}\deleted{(}\handwritten{[}\unsure{\deleted{\illegible{1}}\handwritten{M}}ais on notera que dans le cas d'une catégorie non seulement additi
mais abélienne, il y a une notion plus adéquate du K$_{o}$ , qui devrait corr
pondre à une notion plus adéquate de la suite des K$_{i}$ , inspirée du tapis
de Roos; sans doute c'est en termes de (``vraie'') catégories triangulées

\margin{10}

